% TOP_PROBLEM: {"id": 2302023, "title": "Research Problems in Function Theory — Problem 2.23", "category_id": 8, "category": {"id": 8, "name": "analysis", "display_name": "Analysis", "description": "Limits, continuity, calculus, and function theory.", "slug": "analysis", "order_index": 8, "created_at": "2026-07-31T15:26:25.671Z"}, "status": "open"}
% FIRST_POSTED: null
% ENTRY_KIND: statement_only
% Imported from: batch11.tex; UnsolvedMath ID 2302023
% Compile twice: lualatex 2302023.tex
\documentclass[11pt]{article}
\usepackage{amsmath,amssymb,mathtools}
\usepackage{fontspec}
\usepackage{unicode-math}
\setmainfont{Latin Modern Roman}
\setmathfont{Latin Modern Math}
\usepackage{xurl}
\usepackage[unicode,colorlinks=true,linkcolor=blue,urlcolor=blue,pdftitle={UnsolvedMath Problem 2302023}]{hyperref}
\newcommand{\sref}[2]{\hyperlink{source:#1:#2}{[S#2]}}
\newcommand{\cataloguescope}[1]{\paragraph{Mathematical statement.}#1}
\begin{document}
% UnsolvedMath ID 2302023; source catalogue code AMR-022-2023
\setcounter{section}{2302022}
\section{Research Problems in Function Theory — Problem 2.23}\label{2302023}
{\small\sffamily UnsolvedMath ID 2302023 · Analysis}
\subsection{Definitions and mathematical statement}

\paragraph{Shared notation.}$\mathbb N=\{1,2,\ldots\}$, and $\mathbb Z,\mathbb Q,\mathbb R,\mathbb C$ denote the integers, rationals, reals and complex numbers. Any other conventions are specified locally.


Let $f:\mathbb C\to\mathbb C$ be entire and set $f^{\circ1}=f$, $f^{\circ(n+1)}=f\circ f^{\circ n}$.
A point $z_0$ has exact period $n$ if $f^{\circ n}(z_0)=z_0$ and $f^{\circ k}(z_0)\ne z_0$ for $1\le k<n$.
The source calls such a point a ``fixed point of exact order $n$''.
A straight line means a set $\ell=\{a+tb:t\in\mathbb R\}$ with $b\ne0$.

\paragraph{Statement recorded in the supplied source.}
For every transcendental entire function $f$ and every straight line $\ell$, must there be a periodic point of exact period one or two outside $\ell$? Equivalently, can all points of period at most two lie on one line? \sref{2302023}{2}
Update 2.23 gives a negative answer to this concentration question: for each $n\ge2$ and each line there are infinitely many fixed points of $f^{\circ n}$ off that line, with the later result giving repelling ones. \sref{2302023}{2}
\subsection{Short English statement}
Can all fixed points and all points of exact period two of a transcendental entire function lie on a single straight line?
\subsection{Sources}
\begin{description}
\item[\hypertarget{source:2302023:1}{S1}] ULAM.AI and the UnsolvedMath Contributors, \emph{UnsolvedMath: A Curated Collection of Open Mathematics Problems}, record 2302023 (AMR-022-2023). CC BY 4.0. \url{https://huggingface.co/datasets/ulamai/UnsolvedMath}.
\item[\hypertarget{source:2302023:2}{S2}] W. K. Hayman and E. F. Lingham, \emph{Research Problems in Function Theory} (2018), Chapter 2, Problem 2.23 and Update 2.23. \url{https://arxiv.org/abs/1809.07200}.
\end{description}
\label{end:2302023}
\end{document}
